% cap6_e2c.tex — Seção Cap. 6 E2-C (copiar/adaptar para SLR-Agente)
% Fonte: PathCast docs/thesis/cap6_e2c.tex
% Tabelas: \input a partir de docs/thesis/tables/e2c_*.tex

\section{Experimento cross-paradigma (E2-C)}
\label{sec:e2c}

O experimento E2-C complementa a avaliação interna E2 (RQ1--RQ3) com duas comparações
cross-paradigma sobre o subconjunto de 58 projetos Jira com crosswalk validado para o
audit IJF do MCS\_ML~\cite{...}. A análise primária de throughput restringe-se ao
\emph{evaluation pool} pré-registrado ($n=20$); a análise secundária reporta os 58 projetos.

\subsection{Throughput semanal (RQ4)}
\label{sec:e2c-throughput}

Derivamos a série semanal de conclusões via simulação Markov+Monte Carlo
(B\_TH-Markov), agregada por semana ISO com zeros explícitos, e avaliamos CRPS
(itens/semana) fold-a-fold conforme o protocolo IJF. A métrica legada
\texttt{crps\_weeks} do E2 \emph{não} é comparável ao painel IJF.

% Auto-generated by scripts/generate_e2c_summary_tex.py

\begin{table}[ht]
\centering
\caption{E2-C — throughput semanal (CRPS, itens/semana)}
\label{tab:e2c-throughput-overview}
\begin{tabular}{lrrr}
\toprule
Métrica & Mediana & Min & Max \\
\midrule
B\_TH-Markov & 16.99 & 0.79 & 894.38 \\
Empírico (event log) & 15.86 & 0.76 & 1596 \\
\midrule
\multicolumn{4}{l}{Markov $<$ empírico: 14/51 projetos} \\
\bottomrule
\end{tabular}
\end{table}


\subsection{Lead time em dias (RQ4b)}
\label{sec:e2c-leadtime}

O ground truth canônico é o \emph{span} do event log convertido para dias.
Comparamos PathCast Markov+MC com a baseline B6 (ML-LT refatorado, sem \emph{leakage}
de features temporais pós-abertura).

% Auto-generated by scripts/generate_e2c_summary_tex.py

\begin{table}[ht]
\centering
\caption{E2-C — lead time (CRPS, dias)}
\label{tab:e2c-leadtime-overview}
\begin{tabular}{lrr}
\toprule
Modelo & Mediana CRPS (dias) & $N$ \\
\midrule
PathCast Markov+MC & 122.06 & 51 \\
B6 ML-LT & 126.09 & 7 \\
\midrule
\multicolumn{3}{l}{PathCast $<$ B6: 4/7 projetos (pareado)} \\
\bottomrule
\end{tabular}
\end{table}


\subsection{Inferência estatística}
\label{sec:e2c-stats}

Aplicamos Wilcoxon pareado com correção de Holm no pool evaluation ($n=20$),
separadamente para throughput e lead time.

% Auto-generated by scripts/generate_e2c_summary_tex.py

\begin{table}[ht]
\centering
\caption{E2-C — Wilcoxon pareado (pool evaluation IJF, $n=20$)}
\label{tab:e2c-wilcoxon}
\begin{tabular}{lrrrr}
\toprule
Comparação & $n$ & Vitórias A & $p_{Holm}$ & Sig. \\
\midrule
throughput\_markov\_vs\_empirical & 15 & 2 & 0.9990 & $\times$ \\
throughput\_markov\_crpss\_vs\_empirical\_crpss & 15 & 12 & 0.0103 & \checkmark \\
leadtime\_pathcast\_vs\_b6 & 0 & 0 & --- & $\times$ \\
\bottomrule
\end{tabular}
\end{table}


\subsection{Análise estratificada (sensibilidade)}
\label{sec:e2c-stratified}

O corpus PathCast E2 Jira apresenta \texttt{risk\_level}=HIGH homogêneo (58/58).
Como análise de sensibilidade, estratificamos por quartis de throughput médio do
catálogo MCS\_ML (proxy de intensidade processual; ver Q6 em MCSML-Q-PERTINENCE).

% Auto-generated by scripts/generate_e2c_summary_tex.py

\begin{table}[ht]
\centering
\caption{E2-C — estratificação (pool evaluation IJF)}
\label{tab:e2c-stratified}
\begin{tabular}{llrrrr}
\toprule
Dimensão & Estrato & $n$ & Med. TH Markov & Med. TH Emp. & Vitórias M$<$E \\
\midrule
risk\_level & HIGH & 20 & 30.28 & 22.51 & 2/15 \\
throughput\_tier & Q1\_low & 4 & 1.02 & 0.94 & 0/2 \\
throughput\_tier & Q2 & 5 & 9.47 & 7.01 & 1/4 \\
throughput\_tier & Q3 & 4 & 36.18 & 24.05 & 0/4 \\
throughput\_tier & Q4\_high & 7 & 68.04 & 72.63 & 1/5 \\
\bottomrule
\end{tabular}
\\\textit{\small Uniform risk\_level=HIGH in PathCast E2 corpus (58/58); use throughput\_tier as secondary SPC proxy.}
\end{table}


\subsection{Detalhe por projeto}
\label{sec:e2c-per-project}

% Requires \usepackage{longtable}
% Auto-generated by scripts/generate_e2c_summary_tex.py

% Requires: \usepackage{longtable}

\begin{longtable}{lrrrrrr}
\caption{E2-C — métricas por projeto Jira}\\
\label{tab:e2c-per-project}\\
\toprule
Projeto & $n_{test}$ & TH Markov & TH Emp. & LT PC (d) & B6 (d) & IJF EB \\
\midrule
\endfirsthead
\toprule
Projeto & $n_{test}$ & TH Markov & TH Emp. & LT PC (d) & B6 (d) & IJF EB \\
\midrule
\endhead
\bottomrule
\endlastfoot
\multicolumn{7}{l}{\textit{\small apache}} \\
apache/FLINK & 2240 & 34.92 & 17.57 & 72.25 & --- & --- \\
apache/AIRFLOW & 283 & 2.36 & 2.59 & 94.93 & --- & --- \\
apache/KAFKA & 1176 & 6.52 & 2.64 & 114.65 & --- & 9.48 \\
apache/COMMONS-LANG & 203 & 0.79 & 0.76 & 205.16 & --- & 2.67 \\
apache/MAVEN & 373 & 1.28 & 1.03 & 339.26 & --- & --- \\
\multicolumn{7}{l}{\textit{\small atlassian}} \\
atlassian/CONFCLOUD & 7185 & 60.05 & 202.02 & 307.03 & --- & --- \\
atlassian/BCLOUD & 5592 & 45.19 & 36.80 & 542.20 & --- & --- \\
atlassian/JRACLOUD & 7619 & 118.95 & 112.42 & 587.92 & --- & 39.03 \\
atlassian/JRASERVER & 13379 & 109.65 & 107.97 & 1554 & --- & --- \\
atlassian/CONFSERVER & 12578 & 103.31 & 96.58 & 2059 & --- & --- \\
\multicolumn{7}{l}{\textit{\small hyperledger}} \\
hyperledger/FABN & 481 & 4.66 & 4.12 & 54.11 & 126.09 & --- \\
hyperledger/IS & 430 & 4.44 & 4.76 & 66.05 & --- & --- \\
hyperledger/INDY & 599 & 9.29 & 9.07 & 85.73 & --- & --- \\
hyperledger/FAB & 4043 & 28.30 & 24.86 & 117.65 & --- & 38.52 \\
hyperledger/STL & 479 & 3.95 & 5.43 & 320.46 & --- & --- \\
\multicolumn{7}{l}{\textit{\small inteldaos}} \\
inteldaos/DAOS & 2366 & 43.33 & 30.90 & 31.52 & --- & --- \\
inteldaos/CART & 250 & 1.68 & 1.38 & 60.37 & --- & --- \\
\multicolumn{7}{l}{\textit{\small jfrog}} \\
jfrog/HAP & 228 & 1.37 & 1.20 & 25.30 & --- & --- \\
jfrog/RTFACT & 2705 & 50.07 & 49.30 & 76.75 & --- & --- \\
\multicolumn{7}{l}{\textit{\small jira-ecosystem}} \\
jira-ecosystem/PL & 867 & 3.06 & 6.27 & 70.95 & 73.76 & --- \\
jira-ecosystem/AMKT & 1354 & 3.07 & 4.35 & 171.76 & 155.82 & --- \\
jira-ecosystem/UPM & 1128 & 5.77 & 5.79 & 1606 & --- & --- \\
jira-ecosystem/STRM & 688 & 9.05 & 1596 & 2481 & --- & --- \\
jira-ecosystem/AUI & 1255 & 7.98 & 8.07 & 5955 & --- & 8.88 \\
\multicolumn{7}{l}{\textit{\small mariadb}} \\
mariadb/CONJ & 222 & 1.26 & 1.11 & 19.49 & --- & 2.22 \\
mariadb/MDEV & 5228 & 46.75 & 15.86 & 73.47 & --- & --- \\
mariadb/MXS & 951 & 7.61 & 4.07 & 180.39 & --- & --- \\
mariadb/MCOL & 768 & 10.96 & 5.94 & 208.74 & --- & 6.29 \\
\multicolumn{7}{l}{\textit{\small mindville}} \\
mindville/ICS & 281 & 1.90 & 1.81 & 53.00 & --- & --- \\
\multicolumn{7}{l}{\textit{\small mojang}} \\
mojang/BDS & 4119 & 64.47 & 33.90 & 24.32 & --- & 75.25 \\
mojang/MCCE & 3533 & 68.04 & 72.63 & 48.23 & 56.49 & 43.65 \\
mojang/MC & 64016 & 709.84 & 306.84 & 119.50 & --- & --- \\
mojang/MCL & 5201 & 71.06 & 39.16 & 122.11 & --- & 25.65 \\
mojang/MCPE & 43229 & 894.38 & 602.61 & 16481 & --- & 249.03 \\
\multicolumn{7}{l}{\textit{\small mongodb}} \\
mongodb/WT & 1920 & 25.12 & 21.27 & 50.18 & --- & 16.07 \\
mongodb/JAVA & 1213 & 7.85 & 4.22 & 50.67 & --- & --- \\
mongodb/EVG & 3635 & 42.08 & 13.90 & 55.18 & --- & 20.25 \\
mongodb/DOCS & 4040 & 15.43 & 7.68 & 78.89 & --- & --- \\
mongodb/SERVER & 17374 & 177.87 & 98.46 & 82.57 & --- & --- \\
\multicolumn{7}{l}{\textit{\small qt}} \\
qt/QT3DS & 1139 & 6.76 & 11.06 & 139.00 & --- & --- \\
qt/QDS & 1095 & 31.92 & 22.51 & 241.15 & --- & 12.51 \\
\multicolumn{7}{l}{\textit{\small redhat}} \\
redhat/ENTESB & 4037 & 34.10 & 18.63 & 79.73 & 119.80 & --- \\
redhat/KEYCLOAK & 4533 & 55.59 & 38.50 & 122.06 & --- & --- \\
redhat/JBEAP & 4946 & 33.74 & 19.43 & 205.87 & --- & --- \\
redhat/WFLY & 3612 & 30.28 & 26.82 & 229.85 & --- & 13.49 \\
redhat/JBIDE & 7877 & 33.86 & 26.91 & 2029 & --- & --- \\
\multicolumn{7}{l}{\textit{\small secondlife}} \\
secondlife/STORM & 460 & 1.01 & 2.68 & 262.89 & --- & --- \\
\multicolumn{7}{l}{\textit{\small sonatype}} \\
sonatype/MVNCENTRAL & 493 & 13.63 & 1134 & 1041 & 976.75 & --- \\
sonatype/OSSRH & 22486 & 702.53 & 635.11 & 1320 & 1083 & --- \\
\multicolumn{7}{l}{\textit{\small spring-projects}} \\
spring-projects/SPRING-BOOT & 367 & 2.29 & 3.94 & 39.50 & --- & --- \\
spring-projects/SPRING-FRAMEWORK & 3026 & 16.99 & 15.04 & 122.50 & --- & --- \\
\end{longtable}


\subsection{Limitações}
\label{sec:e2c-limitations}

\begin{itemize}
  \item Unidades distintas: throughput (itens/semana) vs lead time (dias).
  \item Pool evaluation ($n=20$) vs corpus completo ($n=58$).
  \item B6 depende de \texttt{issue\_features.parquet} local (MCS\_ML\_LT).
  \item Reprodutibilidade: bundle Zenodo v0.8.0 (DOI pendente).
\end{itemize}
